\chapter{Two Stages}
\label{ch:two-stage}

\zcref[S]{ex:single-damped} shows the price of meeting the requirement with
one damped stage: about \qty{1}{\milli\henry} of inductance. Two stages get the
same attenuation from far smaller parts, because attenuation grows as $f^4$
instead of $f^2$. The catch is that cascaded stages do not resonate at their
isolated frequencies.

\section{Transfer function}

\begin{derivation}
	Take the lossless ladder with an ideal source and no choke, so
	$Z_{up,0} = sL_1$. The divider at $C_1$ is
	\[
		\frac{Z_{C1}}{Z_{C1} + sL_1} = \frac{1}{1 + s^2 L_1C_1}.
	\]
	Looking upstream from $L_2$, stage 1 is
	\[
		Z_{up,1} = sL_1 \parallel \frac{1}{sC_1} = \frac{sL_1}{1 + s^2L_1C_1}.
	\]
	The divider at $C_2$, after multiplying through by $sC_2$:
	\[
		\frac{1/(sC_2)}{1/(sC_2) + sL_2 + Z_{up,1}}
		= \frac{1}{1 + s^2L_2C_2 + \dfrac{s^2L_1C_2}{1 + s^2L_1C_1}}.
	\]
	Multiply the two dividers. The $(1 + s^2L_1C_1)$ factors combine:
	\[
		H(s) = \frac{1}{(1 + s^2L_1C_1)(1 + s^2L_2C_2) + s^2L_1C_2}.
	\]
\end{derivation}

Expanding,
\begin{equation}\label{eq:H-two-stage}
	H(s) = \frac{1}{1 + s^2 b + s^4 a}, \qquad
	a = L_1L_2C_1C_2, \qquad
	b = L_1C_1 + L_2C_2 + \underbrace{L_1C_2}_{\text{interaction}}.
\end{equation}

\begin{idea}
	If the stages did not load each other, $H$ would just be
	$(1 + s^2L_1C_1)(1 + s^2L_2C_2)$ upside down. The extra $L_1C_2$ term comes
	entirely from stage 2 seeing stage 1's impedance $Z_{up,1}$. It is small only
	when $C_2 \ll C_1$. Otherwise the isolated-stage picture fails, and you must
	design damping for the coupled modes.
\end{idea}

\section{Coupled modes}

With $s^2 = -\omega^2$, the poles satisfy a quadratic in $\omega^2$:
\begin{equation}\label{eq:modes}
	a\,\omega^4 - b\,\omega^2 + 1 = 0
	\quad\Rightarrow\quad
	\omega^2 = \frac{b \pm \sqrt{b^2 - 4a}}{2a}.
\end{equation}

\begin{derivation}
	The product of the roots of a quadratic is its constant term over its
	leading coefficient:
	\[
		\omega_{lo}^2\,\omega_{hi}^2 = \frac{1}{a} = \frac{1}{L_1C_1}\cdot\frac{1}{L_2C_2} = \omega_{01}^2\,\omega_{02}^2
		\quad\Rightarrow\quad
		f_{lo}\,f_{hi} = f_{01}\,f_{02}.
	\]
	Coupling pushes the modes apart, but their product never moves.
\end{derivation}

\subsection{Estimating the modes on sight}

\begin{itemize}
	\item \textbf{Low mode.} At low frequency $L_2$ is nearly a short, so $C_1$
	      and $C_2$ sit in parallel and charge together through $L_1 + L_2$.
	\item \textbf{High mode.} At high frequency $L_1$ is nearly open, so $L_2$
	      rings in a loop with $C_1$ and $C_2$ in series.
\end{itemize}
\begin{equation}\label{eq:mode-estimates}
	f_{lo} \approx \frac{1}{2\pi\sqrt{(L_1 + L_2)(C_1 + C_2)}}, \qquad
	f_{hi} \approx \frac{1}{2\pi\sqrt{L_2 C_s}}, \qquad
	C_s = \frac{C_1C_2}{C_1 + C_2}.
\end{equation}

These mode shapes are what tell you where to put damping (\zcref{sec:damp-ladder}).

\section{Output impedance null}

Between the two peaks, $\Zo$ dips toward zero.

\begin{derivation}
	$\Zo = Z_{C2} \parallel (sL_2 + Z_{up,1})$ goes to zero when the $L_2$
	branch looks like a short:
	\[
		sL_2 + \frac{sL_1}{1 + s^2L_1C_1} = 0
		\quad\Rightarrow\quad
		L_2(1 - \omega^2L_1C_1) + L_1 = 0
		\quad\Rightarrow\quad
		\omega_{null}^2 = \frac{L_1 + L_2}{L_1L_2C_1}.
	\]
	This is a series resonance of $L_2$ against stage 1 (which is capacitive
	above $f_{01}$), shorting $C_2$.
\end{derivation}

\section{High-frequency asymptote}
\label{sec:hf-asymptote}

Well above both modes the $s^4$ term dominates \zcref{eq:H-two-stage}:
\begin{equation}\label{eq:hf-two-stage}
	\abs{H} \approx \frac{1}{\omega^4 L_1L_2C_1C_2} = \left(\frac{f_{01}f_{02}}{f^2}\right)^2,
	\qquad A_{dB} = 40\log_{10}\frac{f^2}{f_{01}f_{02}}.
\end{equation}

The interaction term lives only in $b$, so this asymptote is exact however
strongly the stages interact.

\section{The candidate design}
\label{sec:candidate}

The rest of these notes use this two-stage design, laid out as in
\zcref{fig:dm-circuit}.

\begin{center}
	\begin{tabular}{@{}lll@{}}
		\toprule
		Element                  & Value                                                            & Parasitic                     \\
		\midrule
		CM choke                 & \qty{1}{\milli\henry} CM per winding, \qty{2}{\micro\henry} DM leakage (loop total) & \qty{0.03}{\ohm} per winding  \\
		$L_1$ (source side)      & \qty{39}{\micro\henry}                                           & \qty{0.08}{\ohm} DCR          \\
		Damping leg across $L_1$ & $L_d = nL_1$ in series with $R_d$                                & \qty{0.05}{\ohm} DCR in $L_d$ \\
		$C_1$                    & \qty{2.2}{\micro\farad} effective                                & \qty{10}{\milli\ohm} ESR      \\
		$L_2$ (converter side)   & \qty{4.7}{\micro\henry}                                          & \qty{0.02}{\ohm} DCR          \\
		$C_2$                    & \qty{4.7}{\micro\farad} effective                                & \qty{5}{\milli\ohm} ESR       \\
		\bottomrule
	\end{tabular}
\end{center}

\begin{examplebox}[label=ex:coupled]{where the candidate really resonates}
	Strip the candidate to its four main parts: no parasitics, no damping, no
	choke, ideal source.

	\textbf{Isolated stages.}
	\[
		f_{01} = \frac{1}{2\pi\sqrt{(39)(2.2)}\,\mu\text{s}} = \qty{17.2}{\kHz}, \quad R_{01} = \qty{4.21}{\ohm};
		\qquad
		f_{02} = \qty{33.9}{\kHz}, \quad R_{02} = \qty{1.00}{\ohm}.
	\]

	\textbf{Interaction.} The three terms of $b$ are
	\[
		L_1C_1 = \qty{8.58e-11}{\second\squared}, \qquad
		L_2C_2 = \qty{2.21e-11}{\second\squared}, \qquad
		L_1C_2 = \qty{1.83e-10}{\second\squared}.
	\]
	The interaction term is the \emph{largest} of the three, so the isolated
	picture will be badly wrong. With $a = \qty{1.895e-21}{\second\tothe{4}}$,
	\zcref{eq:modes} gives
	\[
		f_{lo} = \qty{9.44}{\kHz}, \qquad f_{hi} = \qty{61.7}{\kHz}.
	\]
	The modes sit well outside the isolated pair, nearly a factor of 2 each way.
	As a check, $f_{lo}f_{hi} = \qty{5.82e8}{\hertz\squared} = f_{01}f_{02}$.

	\textbf{On sight.} \zcref[S]{eq:mode-estimates} gives \qty{9.17}{\kHz} from
	\qty{43.7}{\micro\henry} against \qty{6.9}{\micro\farad}, and \qty{60.0}{\kHz}
	from $L_2$ against $C_s = \qty{1.50}{\micro\farad}$. Both within 3\%.

	\textbf{Null.} $\Zo$ dips to zero at
	$\sqrt{(L_1+L_2)/(L_1L_2C_1)}/2\pi = \qty{52.4}{\kHz}$, between the two modes.

	\textbf{Attenuation.} At \qty{400}{\kHz} the asymptote gives $-97.6$~dB and
	the exact value is $-97.4$~dB. That is 17~dB more than the 80.4~dB
	requirement, which leaves room for the damping and choke leakage to come.
\end{examplebox}

% -------------------------------------------------------------------

\section{Damping a coupled ladder}
\label{sec:damp-ladder}

The candidate damps stage 1 only: CM choke, $L_1$ with an $R_d$--$L_d$
leg, $C_1$, then an undamped $L_2$--$C_2$. Two questions decide whether that
works. Which mode does the leg damp, and can the undamped stage go without damping?

\subsection{Damp the coupled mode, not the isolated stage}

\begin{idea}
	A leg across $L_1$ damps the mode whose current flows through $L_1$:
	the low mode. In that mode $L_2$ is nearly a short, so $C_1$ and $C_2$
	charge together through $L_1$. The $R_0$ that belongs in the damping
	formulas uses the damped inductor and every capacitance the mode charges:
	\[
		R_{0,low} = \sqrt{\frac{L_1}{C_1 + C_2}}.
	\]
\end{idea}

The closed form reads low when inductance sits outside the damped one
($L_{lk}$ and $L_2$ here), since that inductance is part of the mode but
does not see the leg.

High-frequency attenuation with the leg and the choke leakage uses the
inductance left in series above the leg's corner:
\begin{equation}\label{eq:l1-eff}
	L_{1,eff} = L_{lk} + \frac{nL_1}{1+n}, \qquad
	\abs{H} \approx \frac{1}{\omega^4 L_{1,eff}\,C_1 L_2 C_2}.
\end{equation}

\subsection{Is the undamped stage OK?}
\label{sec:undamped-stage}

\begin{derivation}
	In the high mode $L_1$ is effectively open, so from the $C_2$ terminals
	\[
		\Zo \approx Z_{C2} \parallel Z_b, \qquad
		Z_b = R_b + j\left(\omega L_2 - \frac{1}{\omega C_1}\right),
	\]
	where $Z_b$ is the $L_2$--$C_1$ branch. At the loop resonance the reactances
	around the loop cancel, so $Z_{C2} + Z_b = \Rloop$ and the branch reactance
	equals $+X_{C2}$. Then, as in \zcref{sec:zo-peak},
	\[
		\abs{\Zo}_{peak,hi} \approx \frac{(-jX_{C2})(+jX_{C2})}{\Rloop} = \frac{X_{C2}^2}{\Rloop}.
	\]
	With $\omega^2 = 1/(L_2C_s)$, $X_{C2}^2 = 1/(\omega^2C_2^2)$:
	\begin{equation}\label{eq:zo-hi}
		\abs{\Zo}_{peak,hi} \approx \frac{L_2 C_s}{C_2^2\,\Rloop}, \qquad
		\Rloop = R_{L2} + \ESR_1 + \ESR_2.
	\end{equation}
\end{derivation}

Compare it with $\abs{\Zin}/3.16$ for a 10~dB margin. Watch the scaling:
halving $C_2$ roughly triples it ($C_s$ falls too), and halving $\Rloop$ doubles it.

\begin{figure}[ht]
	\centering
	\begin{tikzpicture}
		\begin{loglogaxis}[
				width=0.9\textwidth, height=6.5cm,
				xmin=3e3, xmax=3e5, ymin=0.05, ymax=80,
				xtick={3e3,1e4,3e4,1e5,3e5}, xticklabels={3k,10k,30k,100k,300k},
				legend style={font=\small},
				xlabel={Frequency (Hz)}, ylabel={$\abs{\Zo}$ ($\Omega$)},
				legend pos=north east, legend cell align=left,
				grid=both, grid style={black!8},
			]
			\addplot[black, thick] table[x=f, y=Zo_undamped] {figures/two-stage-zo.dat};
			\addplot[blue, thick] table[x=f, y=Zo_damped] {figures/two-stage-zo.dat};
			\addplot[red, dashed, thick, domain=3e3:3e5] {26.9};
			\draw[gray, densely dotted] (axis cs:17.2e3,0.05) -- (axis cs:17.2e3,80)
			node[pos=0.06, right, font=\scriptsize] {$f_{01}$};
			\draw[gray, densely dotted] (axis cs:33.9e3,0.05) -- (axis cs:33.9e3,80)
			node[pos=0.06, right, font=\scriptsize] {$f_{02}$};
			\legend{undamped, {$n=1$, $R_d=\qty{3.4}{\ohm}$}, {$\abs{\Zin}/3.16$}}
		\end{loglogaxis}
	\end{tikzpicture}
	\caption{The candidate design with parasitics and choke leakage. The peaks sit at the
		coupled modes, nowhere near the isolated stage resonances $f_{01}$ and $f_{02}$.
		The leg across $L_1$ flattens the low mode.}
	\label{fig:two-stage-zo}
\end{figure}

\begin{examplebox}[label=ex:candidate-damped]{damping the candidate}
	Now the full candidate: parasitics, the \qty{2}{\micro\henry} choke leakage,
	ideal source.

	\textbf{Undamped.} The sweep shows $\Zo$ peaks of \qty{43.0}{\ohm} at
	\qty{9.2}{\kHz} and \qty{8.1}{\ohm} at \qty{61.6}{\kHz}, right at the coupled
	modes. Against \qty{85.0}{\ohm} that is a 5.9~dB margin, which fails the
	10~dB target. $H$ peaks at $+24.5$~dB at \qty{9.2}{\kHz}.

	\textbf{Which $R_0$.} The low mode is the problem, and in it $L_2$ is nearly a
	short, so $C_1$ and $C_2$ charge together through $L_1$:
	\[
		R_{0,low} = \sqrt{\frac{\qty{39}{\micro\henry}}{\qty{6.9}{\micro\farad}}} = \qty{2.38}{\ohm},
		\qquad \text{not} \quad \sqrt{\frac{39}{2.2}}\,\Omega = \qty{4.21}{\ohm}.
	\]

	\textbf{Sizing the leg} with $n = 1$ ($L_d = \qty{39}{\micro\henry}$):
	\begin{center}
		\small
		\begin{tabular}{@{}llll@{}}
			\toprule
			$R_0$ used                     & $R_d$            & $\Zo$ peaks                                                                & Margin  \\
			\midrule
			\qty{2.38}{\ohm} ($C_1 + C_2$) & \qty{3.45}{\ohm} & \qty{6.63}{\ohm} at \qty{10.9}{\kHz}, \qty{4.85}{\ohm} at \qty{63.1}{\kHz} & 22.2~dB \\
			\qty{4.21}{\ohm} ($C_1$ only)  & \qty{6.10}{\ohm} & \qty{7.84}{\ohm} at \qty{9.9}{\kHz}, \qty{4.09}{\ohm} at \qty{63.0}{\kHz}  & 20.7~dB \\
			Brute-force optimum            & \qty{3.40}{\ohm} & \qty{6.63}{\ohm}, \qty{4.86}{\ohm}                                         & 22.2~dB \\
			\bottomrule
		\end{tabular}
	\end{center}
	The coupled-mode $R_0$ lands on the true optimum. The formula predicts
	$2.38\sqrt{6} = \qty{5.82}{\ohm}$; the real peak is 14\% higher because
	$L_{lk}$ and $L_2$ are part of the mode but outside the damped inductor.
	The design uses $R_d = \qty{3.4}{\ohm}$ from here on.

	\textbf{The undamped stage.} At \qty{61.6}{\kHz}, $X_{C2} = \qty{0.550}{\ohm}$
	and $\Rloop = 0.02 + 0.01 + 0.005 = \qty{0.035}{\ohm}$, so \zcref{eq:zo-hi}
	gives $0.550^2/0.035 = \qty{8.6}{\ohm}$ (\qty{8.1}{\ohm} exact). That is well
	under \qty{26.9}{\ohm}, so stage 2 needs no damping of its own. Its $H$ peak is
	only about $+2.7$~dB, because the stage-1 divider is already attenuating at
	\qty{62}{\kHz}.

	\textbf{Attenuation.} $L_{1,eff} = 2 + 19.5 = \qty{21.5}{\micro\henry}$, and
	\zcref{eq:l1-eff} gives $-92.4$~dB at \qty{400}{\kHz} ($-92.2$~dB exact).
	That clears the 80.4~dB requirement by 11.8~dB. The low-mode $H$ peak is
	down to $+9.7$~dB.

	\textbf{Alternative.} $n = 0.5$ with $R_d = \qty{2.17}{\ohm}$ gives a 25.9~dB
	margin and $-89.0$~dB at \qty{400}{\kHz}. More margin than needed, at a cost of 3~dB.
\end{examplebox}
